Compact stars, including neutron stars and boson stars, are among the densest known objects in the universe. Their extreme densities make them crucial laboratories for studying fundamental physics, including nuclear matter equations of state and strong-field gravity (Sennett et al., 2017, [DOI: 10.1103/PhysRevD.96.024002](https://journals.aps.org/prd/abstract/10.1103/PhysRevD.96.024002)). Neutron stars are the remnants of supernova explosions, while boson stars, though hypothetical, represent solutions to Einstein’s field equations involving self-interacting scalar fields (Nelson et al., 2019, [DOI: 10.1088/1475-7516/2019/07/012](https://iopscience.iop.org/article/10.1088/1475-7516/2019/07/012/meta)). Recent research has introduced two-fluid stellar models, where a normal nuclear matter component coexists with an interacting bosonic dark matter component, significantly affecting the internal structure and gravitational properties of compact stars (Guha & Sen, 2024, [DOI: 10.1103/PhysRevD.109.043038](https://journals.aps.org/prd/abstract/10.1103/PhysRevD.109.043038)).  

The presence of dark matter within compact stars could significantly alter their mass-radius relationships and tidal deformability. Studies suggest that dark matter may form halos around neutron stars, influencing their structural properties and the signals observed in gravitational waves (Pitz & Schaffner-Bielich, 2025, [DOI: 10.1103/PhysRevD.111.043050](https://journals.aps.org/prd/abstract/10.1103/PhysRevD.111.043050)). Additionally, constraints from binary neutron star mergers observed by LIGO-Virgo have been used to limit the mass and interactions of dark matter particles (Guo et al., 2023, [DOI: 10.1103/PhysRevD.107.083002](https://journals.aps.org/prd/abstract/10.1103/PhysRevD.107.083002)). The inclusion of two-fluid models has shown that a bosonic dark matter component can significantly soften the equation of state of neutron stars, reducing their maximum mass and shifting their tidal deformability signatures in a way that could be detected by future gravitational wave observations (Bezares & Sanchis-Gual, 2024, [DOI: 10.48550/arXiv.2406.04901](https://arxiv.org/abs/2406.04901)). However, the extent to which dark matter influences compact star properties remains an open question.  

This study aims to model the mass-radius (M-R) relationship of neutron stars and boson stars, incorporating dark matter fractions. It also seeks to analyze tidal deformability (Love number \( k_2 \) vs. compactness \( C \)) to understand how dark matter modifies the response of stars to tidal forces. Additionally, it investigates two-fluid models, where dark matter interacts with nuclear matter through bosonic scalar fields, affecting compact star stability and observables. The study further intends to constrain the possible dark matter fraction within compact stars using gravitational wave observations. These objectives build upon recent work demonstrating that dark matter interactions can lead to significant changes in the merger dynamics of neutron stars (Bezares et al., 2024, [DOI: 10.1103/PhysRevD.110.023013](https://journals.aps.org/prd/abstract/10.1103/PhysRevD.110.023013)).  

The interplay between dark matter and gravitational wave astrophysics is gaining increasing attention. Observations of neutron star mergers have provided the first constraints on exotic matter compositions, including self-interacting dark matter (Bauswein et al., 2023, [DOI: 10.1103/PhysRevD.107.083002](https://journals.aps.org/prd/abstract/10.1103/PhysRevD.107.083002)). Additionally, numerical simulations have shown that bosonic dark matter cores and two-fluid models could leave unique imprints on gravitational wave signals, distinguishable from standard neutron star interactions (Bezares et al., 2019, [DOI: 10.1103/PhysRevD.100.044049](https://journals.aps.org/prd/abstract/10.1103/PhysRevD.100.044049)). These insights will contribute to upcoming observational campaigns by the next generation of gravitational wave detectors, enhancing our understanding of both compact star physics and dark matter properties. By integrating theoretical modeling, numerical simulations, and observational constraints, this study advances our ability to probe dark matter’s role in extreme astrophysical environments.  